%% Task 2 onwards — include this after report_part1.tex in the final report.tex

%% ============================================================
\section{Task 2: Neural Style Transfer Video Pipeline}
%% ============================================================

\subsection{System Overview}

\begin{figure}[h]
\centering
\begin{tikzpicture}[
  box/.style={rectangle, draw, rounded corners, fill=orange!15,
              minimum width=2.2cm, minimum height=0.8cm, font=\small, align=center},
  decision/.style={diamond, draw, fill=yellow!20, font=\small, aspect=2},
  arr/.style={-Stealth,thick}
]
\node[box] (vid)  {Input\\Video};
\node[box, right=0.7cm of vid]  (dec)  {Frame\\Decoder};
\node[box, right=0.7cm of dec]  (mat)  {Matting\\Model\\(U-Net)};
\node[box, right=0.7cm of mat]  (nst)  {NST\\(VGG19)};
\node[box, right=0.7cm of nst]  (comp) {Compositor};
\node[box, right=0.7cm of comp] (enc)  {Video\\Encoder};

\draw[arr](vid)--(dec);
\draw[arr](dec)--(mat);
\draw[arr](dec.south) -- ++(0,-0.5) -| (nst.south);
\draw[arr](mat)--(nst) node[midway,above,font=\tiny]{$\alpha_t$};
\draw[arr](nst)--(comp) node[midway,above,font=\tiny]{$S_t$};
\draw[arr](comp)--(enc);

\node[box, below=1.2cm of nst, fill=purple!15] (style) {Style\\Image};
\draw[arr](style.north) -- (nst.south);

\node[box, below=1.2cm of comp, fill=gray!15] (prev) {$S_{t-1}$\\(temporal)};
\draw[arr](comp.south) -- ++(0,-0.3) -- (prev.east);
\draw[arr](prev.west) -- ++(-0.5,0) |- (nst.west);

\end{tikzpicture}
\caption{Video stylisation pipeline. For each frame $t$, the matting model produces
alpha matte $\alpha_t$, and NST produces stylised frame $S_t$. Temporal consistency is
enforced by initialising NST optimisation from $S_{t-1}$ rather than the raw content
frame. Three compositing variants are produced per frame.}
\label{fig:pipeline}
\end{figure}

\subsection{Human Matting Model}

\subsubsection{Architecture}

The matting model is a U-Net (Ronneberger et al., 2015) with a symmetric
encoder-decoder structure and skip connections. The encoder applies five stages of
Conv $\to$ Conv $\to$ MaxPool, doubling channels at each stage: 64, 128, 256, 512,
1024 (bottleneck). The decoder mirrors this with bilinear upsampling (or transposed
convolutions) and concatenation of corresponding encoder features via skip connections.
The output is a single-channel sigmoid-activated map representing the alpha matte
$\alpha \in [0, 1]$ per pixel.

Input images are resized to $256 \times 256$ prior to inference. The model is trained
from scratch on the AISegment dataset; no pretrained backbone is used.

\subsubsection{Loss Function}

The combined loss is:
\begin{equation}
  \mathcal{L} = \lambda_1 \mathcal{L}_{\text{L1}} + \lambda_d \mathcal{L}_{\text{Dice}}
\end{equation}
with $\lambda_1 = \lambda_d = 0.5$. The L1 term penalises pixel-level absolute error
in alpha values. The Dice loss is:
\begin{equation}
  \mathcal{L}_{\text{Dice}} = 1 - \frac{2\sum \hat{\alpha}_i \alpha_i}{\sum \hat{\alpha}_i^2 + \sum \alpha_i^2 + \epsilon}
\end{equation}
Dice loss is preferred over binary cross-entropy for this task because it is robust to
the class imbalance between foreground (person pixels) and background, which can vary
substantially across portrait images.

\subsubsection{Training Details}

Training uses the Adam optimiser with learning rate $10^{-4}$ and L2 weight decay
$10^{-4}$. A ReduceLROnPlateau schedule reduces the learning rate by a factor of 0.5
after three epochs without improvement in validation IoU. Early stopping (patience = 7
epochs) is applied. Augmentation consists of random horizontal flip, colour jitter
(brightness $\pm0.2$, contrast $\pm0.2$, saturation $\pm0.2$, hue $\pm0.05$), and
random crop (85--100\% of image area). All random seeds are fixed at 42.

A subset of 5,000/500/500 image-matte pairs is used for train/val/test,
subsampled uniformly at random (seed 42) from the full 34,427-pair dataset.

\subsubsection{Results}

\begin{table}[h]
\centering
\caption{Matting model performance on the AISegment test split (500 images).}
\begin{tabular}{lcc}
\toprule
Metric & Value & Target \\
\midrule
IoU  & \textbf{[INSERT]} & $\geq 0.85$ \\
MAD  & \textbf{[INSERT]} & --- \\
\bottomrule
\end{tabular}
\end{table}

\begin{figure}[h]
\centering
\fbox{\rule{0pt}{4.5cm}\rule{\linewidth}{0pt}}
\caption{Training curves for the matting model. Top left: combined loss. Top right: IoU.
Bottom left: MAD. Bottom right: learning rate schedule. \textbf{[Replace with
outputs/matting\_training\_curves.png]}.}
\label{fig:matting_curves}
\end{figure}

\begin{figure}[h]
\centering
\fbox{\rule{0pt}{4cm}\rule{\linewidth}{0pt}}
\caption{Matting visualisation on five video frames. Each row shows: original frame,
predicted alpha matte, foreground cutout with transparent background.
\textbf{[Replace with outputs/matting\_overlay.png]}.}
\label{fig:matting_overlay}
\end{figure}

\subsection{Neural Style Transfer}

\subsubsection{Method}

Following Gatys, Ecker and Bethge (2015), style transfer is framed as an optimisation
over the pixels of an output image $\mathbf{x}$, minimising:
\begin{equation}
  \mathcal{L}_{\text{total}} = \alpha \mathcal{L}_{\text{content}} +
                                \beta  \mathcal{L}_{\text{style}} +
                                \gamma \mathcal{L}_{\text{TV}}
\end{equation}
The content loss is the squared Euclidean distance between feature maps at layer
\texttt{conv4\_2} of a frozen VGG19:
\begin{equation}
  \mathcal{L}_{\text{content}} = \frac{1}{2}\sum_{i,j}(F_{ij} - P_{ij})^2
\end{equation}
where $F$ and $P$ are the generated and content feature maps respectively.

The style loss uses Gram matrices at five layers
(\texttt{relu1\_1}, \texttt{relu2\_1}, \texttt{relu3\_1}, \texttt{relu4\_1},
\texttt{relu5\_1}). The Gram matrix of feature map $F \in \mathbb{R}^{C \times HW}$ is:
\begin{equation}
  G = \frac{F F^\top}{C \cdot H \cdot W}
\end{equation}
The style loss is the mean squared error between generated and target Gram matrices,
averaged across layers. Total-variation regularisation $\mathcal{L}_{\text{TV}}$ penalises
spatial discontinuities:
\begin{equation}
  \mathcal{L}_{\text{TV}} = \sum_{i,j}\bigl(|x_{i,j+1}-x_{i,j}| + |x_{i+1,j}-x_{i,j}|\bigr)
\end{equation}

\subsubsection{Optimisation}

The optimised image $\mathbf{x}$ is a leaf tensor initialised from the content frame.
L-BFGS with strong Wolfe line search is used (100 iterations per frame after frame 0,
300 for frame 0). Loss weights: $\alpha = 10^5$, $\beta = 10^7$, $\gamma = 1.0$.

Temporal consistency is enforced by initialising frame $t$ from the stylised output of
frame $t-1$, rather than from the raw content frame. This substantially reduces
inter-frame flickering at zero additional parameter cost.

\subsubsection{\texorpdfstring{$\beta/\alpha$}{beta/alpha} Ablation}

\begin{figure}[h]
\centering
\fbox{\rule{0pt}{4cm}\rule{\linewidth}{0pt}}
\caption{$\beta/\alpha$ ablation. From left to right: content image, style image,
output at $\beta = 10^4$, $\beta = 10^6$, $\beta = 10^8$ (with $\alpha$ fixed at $10^5$).
\textbf{[Replace with outputs/beta\_alpha\_ablation.png]}.}
\label{fig:ablation_ba}
\end{figure}

At low $\beta/\alpha$ ratios the output is perceptually similar to the content image.
As the ratio increases, style textures become progressively dominant. At $\beta/\alpha =
10^3$ the compositional structure of the painting transfers but semantic content is
largely preserved. At $\beta/\alpha = 10^5$ both structure and fine texture are imposed,
and content identity is largely lost.

\subsubsection{Layer Ablation}

\begin{figure}[h]
\centering
\fbox{\rule{0pt}{4cm}\rule{\linewidth}{0pt}}
\caption{Layer ablation. Left to right: content, style, shallow-only
(relu1\_1, relu2\_1), deep-only (relu4\_1, relu5\_1), full five-layer result.
\textbf{[Replace with outputs/layer\_ablation.png]}.}
\label{fig:layer_ablation}
\end{figure}

Shallow layers capture grain, edge orientations and colour statistics at fine scales.
Deep layers capture compositional structures and large-scale stroke patterns. The
five-layer combination transfers both, consistent with the original Gatys et al. finding.

\subsection{Video Compositing}

For each frame $F_t$, the matting model produces $\alpha_t \in [0,1]^{H\times W}$,
and NST produces the stylised frame $S_t$. Three variants are composited per frame:
\begin{align}
  O^{\text{bg}}_t   &= \alpha_t \cdot F_t + (1-\alpha_t) \cdot S_t
    \quad \text{(background stylised)} \\
  O^{\text{fg}}_t   &= \alpha_t \cdot S_t + (1-\alpha_t) \cdot F_t
    \quad \text{(subject stylised)} \\
  O^{\text{full}}_t &= S_t
    \quad \text{(full-frame stylised)}
\end{align}
Frames are re-encoded at the original frame rate (30\,fps) using OpenCV.

\begin{figure}[h]
\centering
\fbox{\rule{0pt}{5cm}\rule{\linewidth}{0pt}}
\caption{Sample frames from the three stylised video variants.
Row 1: background stylised. Row 2: subject stylised. Row 3: full-frame stylised.
\textbf{[Replace with extracted frames from output videos]}.}
\label{fig:video_frames}
\end{figure}

\subsection{VGG19 Feature Map Visualisation}

\begin{figure}[h]
\centering
\fbox{\rule{0pt}{5cm}\rule{\linewidth}{0pt}}
\caption{Eight channels from \texttt{relu1\_1} (shallow) and \texttt{relu5\_1} (deep),
shown for one video frame (top) and one seed image from Task 1 (bottom).
\textbf{[Replace with outputs/feature\_maps.png]}.}
\label{fig:feature_maps}
\end{figure}

The shallow layer \texttt{relu1\_1} responds to oriented edges and colour gradients ---
in seed images this corresponds to seed boundaries, and in video frames to hair, clothing
edges, and facial contours. The deep layer \texttt{relu5\_1} responds to more complex
patterns: in seed images, clusters of seeds activate coherent regions, mirroring the
spatial counting signal the Task~1 CNN must learn to integrate.

%% ============================================================
\section{Cross-Method Comparison}
%% ============================================================

\begin{table}[h]
\centering
\caption{Unified comparison across all seed-counting methods. Latency is measured on a
single NVIDIA T4 GPU. Training cost is an estimate based on observed Kaggle session times.
Failure cases are from the failure\_cases.json produced by the Assignment 2 evaluator.}
\begin{tabular}{lcccccc}
\toprule
Method & MAE & RMSE & Latency (ms/img) & Model size & Train cost & Failures fixed \\
\midrule
Clustering (A1)     & $\sim$18 & $\sim$24 & 12  & 0\,MB      & None         & --- \\
Edge Detection (A2) & $\sim$14 & $\sim$19 & 35  & 0\,MB      & None         & 4/N \\
CNN Model A         & [INSERT] & [INSERT] & 4   & 0.4\,MB    & $\sim$4\,GPU-hr & [X/N] \\
CNN Model B         & [INSERT] & [INSERT] & 11  & 18.5\,MB   & $\sim$6\,GPU-hr & [X/N] \\
\bottomrule
\end{tabular}
\label{tab:comparison}
\end{table}

The clustering and edge-detection methods require no training but are brittle to seed
density and lighting variation. Both CNN models generalise across these conditions through
learned feature representations. Model B's additional parameters do not yield substantial
MAE gains over Model A on this dataset size, suggesting the primary bottleneck is data
quantity rather than model capacity.

%% ============================================================
\section{Cost and Deployment Discussion}
\label{sec:cost}
%% ============================================================

\textbf{CNN seed counter.} Model A inference takes approximately 4\,ms per image on a
T4 GPU and 60--80\,ms on a mid-range CPU, making it viable for real-time field deployment
on an edge device. The model file is 0.4\,MB, suitable for embedded systems. Quantisation
to INT8 is expected to halve latency and memory footprint without significant accuracy
loss.

\textbf{Matting model.} U-Net inference on a $256\times256$ frame takes approximately
15\,ms on GPU. At 30\,fps, GPU inference can keep pace with real-time video if the matting
is the only operation being performed. On CPU, this rises to 400--800\,ms, ruling out
real-time CPU deployment.

\textbf{NST pipeline.} NST as implemented (L-BFGS with 100 iterations) requires
approximately 45--90 seconds per frame on a T4 GPU. This is fundamentally incompatible
with real-time operation. For production use, two alternatives exist: (1) fast style
transfer methods (e.g. Johnson et al., 2016) train a feed-forward network to perform
stylisation in a single forward pass, reducing latency to under 10\,ms per frame; (2)
offline batch processing at production time, with the stylised video stored and replayed.

\textbf{Recommendation for AgriVision.} For the seed counting task, deploy Model A on
a GPU-enabled cloud inference endpoint. The 4\,ms latency supports batch processing of
entire field surveys within minutes. For the video marketing content, use offline NST
with a pre-computed style to produce assets; do not attempt real-time stylisation without
a feed-forward architecture. The matting model can run in real-time on GPU as a
pre-processing step to enable selective compositing.

%% ============================================================
\section{Limitations and Future Work}
%% ============================================================

The seed counting CNN is trained on a fixed image size of $128\times128$ and on images
from a single acquisition setup. Performance on images from different cameras, lighting
rigs, or seed varieties is unknown and likely to degrade. The regression formulation
discards information about seed positions; a detection-based approach (e.g. a FCOS or
CenterNet model) would produce counts as a by-product of localisation and would be more
interpretable in failure cases.

The matting model is trained on portrait images from the AISegment dataset, which may
not match the statistics of the user's self-recorded video in terms of background
complexity, clothing, or motion blur. Domain adaptation or fine-tuning on a small set of
video frames would likely improve matte quality.

NST produces temporally consistent output through warm initialisation but does not
enforce optical-flow-based consistency. High-motion frames may still exhibit colour drift.
A future system should incorporate optical flow warping between consecutive stylised frames.

%% ============================================================
\section{Reflection}
%% ============================================================

The most unexpected finding in Task~1 was the magnitude of the optimizer effect: SGD
outperforming Adam by a factor of five was not anticipated. In hindsight this makes sense
for a small dataset where Adam's adaptive second-moment estimates are noisy, but it was
not the hypothesis going into the ablation. This finding shifted the entire subsequent
strategy, as running 768 exhaustive grid-search configurations would have produced results
dominated by Adam's instability rather than meaningful architectural comparisons.

The regression reformulation of the seed-counting task was a deliberate choice made early
in the project. The classification framing in the assignment specification is reasonable
for small count ranges, but impractical for a count distribution spanning 1 to 150 with
sparse examples at the extremes. This decision simplified training and evaluation and
produced interpretable results with a single metric.

For Task~2, the primary technical difficulty was compatibility between the PyTorch 2.x
optimizer internals and L-BFGS, specifically the requirement for contiguous gradient
tensors and the need to explicitly enable gradient tracking within the closure. These are
not documented clearly in the PyTorch L-BFGS documentation and required systematic
debugging.

If the project were repeated, an early experiment to profile dataset-level statistics
(seed count distribution, image-to-image variance) would have informed the choice between
classification and regression before any model was built. Similarly, a brief compatibility
check of L-BFGS on the target PyTorch version before committing to it as the NST
optimiser would have avoided several hours of debugging.

\bigskip
\begin{flushright}
\textit{Report generated: May 2026}
\end{flushright}

\end{document}
